\pdexercisesfor{\S\ref{sec:stp-consolidation}}{Consolidation compacts the view, never the record}

\begin{pdexercise}[kind=Check,rating=1]{ex:stp-consolidation-f4}
Mem0's update step applies one of \texttt{ADD}, \texttt{UPDATE},
\texttt{DELETE} or \texttt{NOOP} to each candidate fact, and \texttt{DELETE}
removes a stored memory that new information contradicts~\pdcite{chhikara2025mem0}.
Suppose the stored memory is a grant note with $\mathsf{ptr}=\{g_1\}$ and a new
message says the grant was revoked. Which condition of
Design invariant~\ref{inv:stp-ledger-faithful} does \texttt{DELETE} violate, and
what should the update write instead?
\end{pdexercise}
\begin{pdsolution}{ex:stp-consolidation-f4}
It violates F4: a note with nonempty $\mathsf{ptr}$ is never deleted or
rewritten in place. Deleting it erases the record of what the agent believed
when it last acted, so an audit or appeal cannot reconstruct that belief. The
update should append a superseding note with $\mathsf{ptr}=\{g_1,r_1\}$, where
$r_1$ is the revocation event in $L$, and tombstone the old note. F1 adds a
second requirement: if the revocation exists only in the message and not in
$L$, the new note is texture and cannot change the grant's status until the
issuer records the revocation. Mem0's graph variant, which marks a conflicting
relation invalid instead of removing it, is closer to F4.
\end{pdsolution}
\begin{pdexercise}[kind=Trace,rating=2]{ex:stp-lunches-trace}
Trace Numbers by Hand~\ref{ex:stp-lunches-only} through the seven steps of the
mediation rule in \S\ref{sec:stp-consolidation}. At which step is the false
authority first stopped? If that step is misconfigured and the widened note is
admitted, which later step stops the 6{,}000 USD order, and what does the
consolidation receipt let an auditor find?
\end{pdexercise}
\begin{pdsolution}{ex:stp-lunches-trace}
Step~1 labels the ERP content with an external source class taken from its
channel. Step~2 does not fire, because the ERP record does not claim to be a
grant; it is texture. Step~3 stops the false authority: the batch would set the
note with $\mathsf{ptr}=\{g_1,g_2\}$ to scope $\{\text{lunch},\text{reception}\}$,
while $L$ gives $S_L=\{\text{lunch}\}$, so F1 fails and the batch is rejected
whole. If the checker is misconfigured and admits the batch, step~7 still stops
the order: admission evaluates $S_L$, reception $\notin S_L$, and the request is
denied. The receipt records the digests before and after, the consolidator and
its engine, the policy version and the verdict. In the misconfigured run it
shows an admitted batch that superseded the $\{g_1,g_2\}$ note, which is where
an auditor finds the fault.
\end{pdsolution}
\begin{pdexercise}[kind=Open,rating=3]{ex:stp-bounded-coverage}
A role carries $k=200$ open obligations. Each has a deadline drawn from $256$
working slots and a counterparty drawn from $16$ principals, chosen
independently. The memory budget is $256$ bytes. Use
Theorem~\ref{thm:stp-coverage-ledger} to show that memory alone cannot satisfy
F2. Then propose the cheapest serving scheme that does, and say what it must
read at boot and what it must do when that read fails.
\end{pdexercise}
\begin{pdsolution}{ex:stp-bounded-coverage}
Each obligation carries $b=\log_2 256+\log_2 16=8+4=12$ independent bits, so
the role needs $kb=200\cdot 12=2{,}400$ bits. The budget is
$256\cdot 8=2{,}048$ bits, and $2{,}048<2{,}400$, so by
Theorem~\ref{thm:stp-coverage-ledger} some assignment is served with wrong terms.
The cheapest scheme stores no terms in memory. At boot the harness reads the
role's open set from $\mathrm{Owed}(L)$ and serves each obligation with its full
terms; obligations created during the session are served from their ledger
events as they are recorded. Memory may hold pointers or none. If the ledger
read fails, the successor must not act within the scope of any obligation it
cannot see: fail closed for those actions, and report the failed read. Open
part: which actions lie outside every obligation's scope, and so stay safe while
the ledger is unreachable?
\end{pdsolution}
